\chapter{Fundamentos de las aplicaciones móviles inteligentes}
\label{chap:fundamentos-moviles}

\insertminitoc
\parindent1.5em

\section{Evolución y arquitectura de las aplicaciones móviles}

La arquitectura de las aplicaciones móviles ha experimentado una evolución drástica, abandonando las estructuras monolíticas tradicionales para adoptar sistemas complejos de múltiples capas. Este cambio de paradigma busca equilibrar la creciente demanda de funcionalidades avanzadas con las severas limitaciones de hardware inherentes a los dispositivos móviles, especialmente aquellos de gama media y baja que predominan en el mercado infantil y adolescente. Según \cite{koulaxidis_improving_2022}, el proceso de ingeniería de software para móviles requiere considerar meticulosamente las restricciones de memoria \gls{ram}, ciclos de \gls{cpu} y disipación térmica. Una arquitectura robusta no solo se encarga de renderizar la interfaz de usuario con fluidez, sino que debe gestionar de forma eficiente la persistencia de datos masivos y la comunicación en tiempo real con sensores internos, garantizando que la aplicación sea estable bajo diversas y extremas condiciones de carga de trabajo.\\

Para abordar esta complejidad, la ingeniería de software moderna propone enfoques como el Desarrollo Dirigido por Modelos (MDD) y la implementación de la Arquitectura Limpia (\textit{Clean Architecture}). Marcos de trabajo recientes como \textit{AppCraft} han demostrado que es posible generar aplicaciones móviles seguras y modulares separando estrictamente la lógica de dominio, la capa de datos y la interfaz de presentación \cite{alwakeel_appcraft_2025}. Esta modularidad es vital en aplicaciones de supervisión infantil, ya que permite a los desarrolladores actualizar los algoritmos de detección de riesgos o los esquemas de bases de datos sin necesidad de reescribir la lógica de la interfaz. Asimismo, una arquitectura desacoplada facilita la realización de pruebas unitarias exhaustivas, asegurando que los fallos en un componente periférico no provoquen el colapso del servicio de monitoreo crítico.\\



En el ecosistema específico de Android, la arquitectura se fundamenta en la interacción de componentes esenciales como \textit{Activities}, \textit{Services}, \textit{Broadcast Receivers} y \textit{Content Providers}. La comunicación entre estos elementos se realiza a través de un mecanismo de Comunicación Inter-Procesos (IPC) basado en \textit{Binder}. Sin embargo, investigaciones técnicas de ciberseguridad han señalado que esta flexibilidad arquitectónica introduce vectores de ataque significativos si los permisos y las interfaces IPC no se gestionan con políticas de seguridad estrictas \cite{rojas_technical_2025, nguyen_security_2021}. Para una aplicación de supervisión inteligente, la arquitectura debe ser particularmente resiliente; el motor de análisis que captura la pantalla debe operar en hilos de ejecución independientes (\textit{Worker Threads}) para no bloquear el hilo principal (\gls{ui} Thread) del sistema operativo, evitando los temidos errores de "La aplicación no responde" (ANR).\\

La tendencia contemporánea más disruptiva en la arquitectura móvil es el desplazamiento del centro de cómputo desde la nube hacia el propio dispositivo. La integración de sistemas inteligentes distribuidos permite que los teléfonos móviles dejen de ser simples terminales tontos de visualización para convertirse en nodos de procesamiento activo y autónomo \cite{oladimeji_smart_2023}. Este enfoque arquitectónico, denominado "Edge-centric" (centrado en el borde), es el pilar fundamental para la protección de la privacidad infantil. Al filtrar y analizar la información sensible directamente en el hardware local antes de cualquier intento de sincronización, se reduce drásticamente la superficie de exposición a ciberataques. A continuación, se presenta una tabla que resume las diferencias entre los paradigmas:\\

\begin{table}[h!]
\centering
\caption{Comparativa de paradigmas arquitectónicos en aplicaciones móviles}
\begin{tabular}{|p{3.5cm}|p{4.5cm}|p{4.5cm}|p{2cm}|}
\hline
\textbf{Característica} & \textbf{Arquitectura Basada en Nube} & \textbf{Arquitectura Local (Edge)} & \textbf{Referencia} \\ \hline
Lugar de Procesamiento & Servidores Remotos centralizados. & Procesador del Dispositivo (\gls{npu}/GPU). & \cite{oladimeji_smart_2023} \\ \hline
Latencia de Respuesta & Alta y variable (depende de la red). & Mínima (tiempo real, $\sim$milisegundos). & \cite{koulaxidis_improving_2022} \\ \hline
Privacidad del Dato & Vulnerable durante el tránsito. & Protegida (aislada en el dispositivo). & \cite{feal_angel_2020} \\ \hline
Consumo de Batería & Alto por transmisión de radio (Wi-Fi/5G). & Optimizado mediante aceleradores locales. & \cite{ali_efficient_2021} \\ \hline
\end{tabular}
\end{table}

\section{Gestión de recursos, servicios en segundo plano y rendimiento del \gls{os}}

El rendimiento de una aplicación móvil de seguridad está intrínsecamente ligado a su capacidad de persistencia, es decir, su habilidad para mantener servicios en segundo plano (\textit{background services}) de forma ininterrumpida sin degradar la experiencia del usuario ni agotar la batería. La optimización del rendimiento es una cuestión de supervivencia frente a los agresivos gestores de memoria de sistemas operativos como Android, que destruyen procesos rutinariamente para liberar memoria \gls{ram} (\textit{Out Of Memory Killer}) \cite{koulaxidis_improving_2022}. Para una aplicación de monitoreo, esto obliga al uso de Servicios de Primer Plano (\textit{Foreground Services}) anclados a notificaciones persistentes, lo cual indica al sistema operativo que la tarea es crítica para el usuario. No obstante, esto impone una responsabilidad ética y técnica al desarrollador de minimizar la huella de memoria (Memory Footprint).\\

Para contrarrestar el agotamiento de la batería, la ingeniería de software móvil ha adoptado técnicas de "Escalado de Calidad Adaptativo" (\textit{Adaptive Quality Scaling}). Según estudios de \cite{ali_efficient_2021}, es posible reducir el consumo energético hasta un 30\% ajustando dinámicamente la frecuencia de muestreo de los sensores y la resolución de las capturas de pantalla en función de la carga térmica de la batería y el nivel de energía restante. Si el teléfono detecta que está por debajo del 15 porciento de batería, el servicio de monitoreo inteligente debe reducir su tasa de fotogramas por segundo (FPS) en el análisis visual o priorizar únicamente el análisis de texto (que es computacionalmente más barato), garantizando así que el dispositivo siga operativo para llamadas de emergencia.\\



El desafío definitivo en el rendimiento móvil radica en la extrema fragmentación del ecosistema Android. Fabricantes como Xiaomi, Samsung y Huawei implementan capas de personalización con políticas de ahorro de energía draconianas que "matan" incluso los servicios de primer plano si sospechan un consumo anómalo. La literatura técnica sugiere que la robustez de una aplicación depende de mecanismos de auto-sanación (\textit{Watchdogs}) y de una interfaz que guíe educativamente al tutor para que otorgue los permisos de "Ignorar optimización de batería" \cite{rojas_technical_2025, ayyash_cybersecurity_2024}. Esto asegura que la prevención de riesgos y el registro en la bitácora de evidencias operen con la fiabilidad de un servicio de sistema continuo.\\

\section{Inteligencia Artificial y Computación en el Borde (Edge AI)}

La evolución de las aplicaciones móviles hacia verdaderas "herramientas inteligentes" es consecuencia directa del auge de la Computación en el Borde (\textit{Edge Computing} y \textit{Edge AI}). Este paradigma postula que el procesamiento intensivo de datos y la ejecución de modelos de aprendizaje automático (Machine Learning) deben ocurrir en los extremos de la red —es decir, en el propio teléfono inteligente— en lugar de delegarse a enormes centros de datos \cite{sahu_edge_2025}. Para la protección infantil, el Edge AI representa un salto cualitativo: permite que el teléfono del menor identifique de forma autónoma, y en milisegundos, imágenes de autolesiones, lenguaje de odio o depredadores sexuales utilizando redes neuronales convolucionales (\gls{cnn}) entrenadas previamente e integradas en el código fuente de la app.\\



La ejecución local de algoritmos de inteligencia artificial requiere técnicas avanzadas de compresión de modelos. Redes inmensas deben ser reducidas mediante procesos de "cuantización" (reduciendo la precisión de los cálculos de 32 bits a 8 bits) y "poda" (\textit{pruning}, eliminando conexiones neuronales redundantes) sin perder precisión predictiva \cite{surianarayanan_survey_2023}. Al utilizar frameworks como TensorFlow Lite o PyTorch Mobile, la aplicación es capaz de enrutar los cálculos tensoriales hacia los aceleradores de hardware específicos del teléfono (\gls{npu} o GPU), liberando a la \gls{cpu} principal. Este nivel de optimización es idéntico al que utilizan las aplicaciones de Realidad Aumentada Móvil (MAR) para reconocer objetos físicos en tiempo real sin vaciar la batería del usuario \cite{cao_mobile_2023}.\\

La autonomía que proporciona la \gls{ia} en el borde tiene un subproducto invaluable: la protección absoluta de la privacidad por diseño. Al no existir la necesidad técnica de transmitir capturas de pantalla de la actividad del menor a través de Internet para su análisis, se erradica el riesgo de interceptación (ataques Man-In-The-Middle) o de almacenamiento no autorizado en servidores de terceros \cite{ali_getting_2023, feal_angel_2020}. Esta arquitectura responde directamente a las demandas éticas de las familias contemporáneas, que se muestran justificadamente reacias a utilizar servicios en la nube que podrían sufrir brechas masivas de seguridad, exponiendo la vida íntima de sus hijos al dominio público.\\

Además, el futuro de estas aplicaciones inteligentes apunta hacia el Aprendizaje Federado (\textit{Federated Learning}). Bajo este esquema, los modelos de \gls{ia} en cada dispositivo móvil aprenden de los nuevos patrones de riesgo que detectan localmente. En lugar de enviar los datos personales del usuario, el teléfono solo envía a la nube "actualizaciones de los pesos matemáticos" de la red neuronal. El servidor central agrega estas actualizaciones anónimas para mejorar el modelo global y luego distribuye la \gls{ia} mejorada a todos los dispositivos \cite{oladimeji_smart_2023}. Esta sinergia entre inteligencia artificial avanzada y preservación de la privacidad marca la frontera tecnológica de la protección infantil del mañana.\\

\section{Ciberseguridad, privacidad por diseño y vulnerabilidades del ecosistema}

El desarrollo de aplicaciones móviles está marcado por una tensión permanente entre la oferta de funcionalidades innovadoras y la protección de los datos personales. Un análisis crítico sobre la privacidad en el ecosistema de aplicaciones de control parental reveló resultados alarmantes: la vasta mayoría de estas herramientas incorpora bibliotecas y rastreadores (\textit{trackers}) ocultos que extraen la geolocalización, hábitos de uso y perfiles de los menores para enviarlos a conglomerados publicitarios sin un consentimiento explícito ni transparente \cite{feal_angel_2020}. Esta "paradoja de la privacidad" subraya que muchas aplicaciones diseñadas teóricamente para proteger a la infancia terminan mercantilizando sus datos, convirtiéndose en el eslabón más débil de la seguridad familiar.\\

Las vulnerabilidades de seguridad trascienden el código de la aplicación y se extienden a las capas de red y sistemas operativos. Investigaciones de \cite{rojas_technical_2025} demostraron que los populares filtros de DNS Protegidos (PDNS), frecuentemente vendidos por los proveedores de internet como soluciones definitivas, son fácilmente eludidos por aplicaciones de VPN de terceros o mediante características integradas en navegadores modernos como el "DNS over HTTPS" (DoH). De manera similar, los fallos de seguridad en entornos interconectados, como el Internet de las Cosas (\gls{iot}) orientado a menores (ej. monitores de bebés), han demostrado que los protocolos de autenticación débiles y la falta de cifrado exponen transmisiones audiovisuales íntimas a la red pública \cite{schmidt_assessing_2023}. Esto evidencia empíricamente que la seguridad perimetral es ineficaz; la protección debe residir directamente en el dispositivo final (\textit{Endpoint Security}).\\



La teoría de la Gestión de la Privacidad en la Comunicación (CPM, por sus siglas en inglés) sugiere que la adopción de tecnologías móviles está fuertemente dictada por cómo el usuario percibe el riesgo de privacidad y cómo evalúa las políticas de control de la aplicación \cite{balapour_mobile_2020}. Si los padres intuyen que una aplicación actúa como un \textit{spyware} que abre "puertas traseras" a corporaciones, la desinstalarán. Para cultivar la confianza, las aplicaciones móviles inteligentes deben adoptar el marco legal e informático de la "Privacidad por Diseño" (Privacy by Design). Esto implica, por ejemplo, que las evidencias de la bitácora de riesgos se cifren utilizando algoritmos de grado militar (ej. AES-256) y que las claves criptográficas se deriven exclusivamente del PIN que solo el padre conoce, asegurando que ni siquiera los desarrolladores de la aplicación puedan leer los registros.\\

Finalmente, las vulnerabilidades a menudo radican en la interfaz de configuración de permisos de Android. Los desarrolladores deben solicitar el acceso a la \gls{api} de Accesibilidad o al Servicio de Proyección de Medios (Screen Capture) explicando de forma inequívoca su propósito. Una auditoría de permisos, sumada a un diseño que no intente ocultar la aplicación al menor —sino que muestre notificaciones claras de que el dispositivo está siendo protegido de forma segura localmente—, alinea la herramienta con los estándares de ciberseguridad exigidos por marcos internacionales de bienestar digital \cite{ayyash_cybersecurity_2024, noauthor_full_2025}.

\section{Desafíos de Usabilidad (HCI) y adopción de tecnologías de seguridad}

El éxito tecnológico de una aplicación móvil no se mide únicamente por la eficiencia de sus algoritmos o la robustez de su arquitectura, sino por su usabilidad en el mundo real. La disciplina de la Interacción Humano-Computadora (HCI) revela que existe una brecha dramática entre las necesidades de seguridad de las familias y su alfabetización digital. Según estudios de \cite{ayyash_cybersecurity_2024}, aunque la ciberseguridad es considerada una prioridad fundamental, la mayoría de los padres abandonan las herramientas de protección debido a interfaces de usuario (\gls{ui}) excesivamente técnicas y difíciles de configurar. Si una aplicación requiere que el tutor entienda conceptos de redes o configure permisos ocultos en submenús del sistema, el diseño ha fracasado en su misión de democratizar la seguridad.\\

El principal problema de usabilidad en los sistemas de monitoreo es la "fatiga de alerta" (\textit{Alert Fatigue}). Los paneles de control tradicionales inundan a los padres con gráficos, listas de sitios web y conteos de horas de pantalla sin ofrecer contexto. Esto genera una carga cognitiva insostenible, provocando que los usuarios ignoren las notificaciones reales. Un diseño centrado en el usuario exige que la Inteligencia Artificial actúe como un filtro inteligente que presente únicamente resúmenes ejecutivos y "momentos clave" procesables \cite{gnanasekaran_usability_2023}. Al traducir los datos técnicos crudos (ej. "intento de acceso a URL maliciosa") en narrativas comprensibles (ej. "Posible contacto con acosador en la aplicación X"), la interfaz facilita que el padre tome decisiones informadas rápidas.\\

Además, el diseño de la aplicación debe concebirse como una herramienta para dos perfiles de usuario radicalmente diferentes: el tutor y el menor. Las investigaciones de diseño participativo demuestran que las interfaces que excluyen la visión del niño generan fricción y desconfianza, resultando en intentos del menor por desinstalar o "hackear" el bloqueo \cite{mcnally_co-designing_2018}. La usabilidad ética sugiere que la aplicación incluya interfaces amigables para el menor, donde él mismo pueda ver (en modo solo lectura) las alertas que se han enviado a sus padres y comprender el motivo de las mismas, fomentando la transparencia y la educación sobre riesgos.\\

Por último, la aplicación móvil inteligente debe trascender el concepto clásico de "tablero de control" (Dashboard) para convertirse en una plataforma narrativa y educativa. Los desarrolladores deben integrar \textit{micro-aprendizajes} dentro de la interfaz, ofreciendo consejos rápidos sobre cómo el padre puede abordar una conversación incómoda si la bitácora detecta contenido relacionado con autolesiones o ciberacoso \cite{ayyash_cybersecurity_2024, stoilova_parental_2024}. Esta integración entre tecnología invisible (Edge AI, optimización de kernel) e interfaces empáticas (HCI, educación narrativa) define el verdadero fundamento de las aplicaciones móviles inteligentes del futuro, posicionándolas no como software de espionaje, sino como mediadores esenciales para la salud familiar.\\